% ============================================
% Chapter: Bill of Materials
% ============================================

\chapter{Bill of Materials}
\label{ch:appendix-a}

Appendix A: Bill of Materials
Table A.1: Bill of Materials --- GreenLab hardware node, as specified in Table 4.2.

Prices are indicative retail unit prices in USD as of mid-2026, based on typical hobbyistelectronics marketplace pricing (Amazon, AliExpress, and comparable distributors).
Actual prices vary by supplier, order volume, and region; verify current pricing with
your chosen supplier before including exact figures as a submitted financial
commitment.

Qty Unit Line
# Component Spec / Model per price total
node (USD) (USD)
ESP32 DevKit V1
Microcontroller $7.50 $7.50
(WROOM-32, WiFi/BT)

Ultrasonic distance
HC-SR04 $2.00 $2.00
sensor

Temperature/humidity
DHT11 $2.20 $2.20
sensor

Light-dependent LDR + comparator
$1.20 $1.20
resistor module breakout

5V single-channel, opto-
Relay module isolated (SRD-05VDC-SL- $2.50 $5.00
C based)

M-F / M-M assortment,
Jumper wires 1 pack $3.50 $3.50
20cm

Breadboard or Half-size prototyping
$3.00 $3.00
perfboard board

5V/2A wall adapter +
USB power supply $4.50 $4.50
micro-USB cable


Enclosure ABS project box, $6.00 $6.00
ventilated,

Qty Unit Line
# Component Spec / Model per price total
node (USD) (USD)
~100×70×40mm

Screws, standoffs, cable
Mounting hardware 1 set $1.50 $1.50
ties

Node
$36.40
subtotal


Per-laboratory deployment cost (one node per lab; IP camera assumed already present as
existing infrastructure, per the Section 1.6 scope boundary): approximately $36--40 per
laboratory, before labor, shipping, and bulk-order discount. Relay, sensor, and
microcontroller unit prices typically drop 15--30% at 10+ unit order quantities from most
suppliers, which is relevant to the Phase 3 pilot-batch expansion proposed in Section 6.5.

Table A.2: Alternative sensor configuration for scaled future deployments.

Presented as a documented option for Phase 4 field-pilot consideration (Section 6.5), not
as a replacement for Table A.1 --- the current design and Table 4.2 specify DHT11 and HCSR04.

Unit
Component Spec / Model price Rationale
(USD)
Roughly 3× the cost of DHT11, but
meaningfully tighter accuracy and wider
DHT22 ±0.5$^\circ$C / ±2--5%
$6.50 operating range --- worth considering if
(AM2302) RH accuracy
Phase 4 field data shows DHT11 accuracy
is a limiting factor.

Would supplement, not replace, the HCPIR motion Passive infrared, SR04: PIR detects heat-signature motion
sensor (HC- adjustable $1.80 (lower false-positive rate for human
SR501) sensitivity/delay presence) while HC-SR04 detects
proximity/distance. Combining both
would strengthen the multimodal

Unit
Component Spec / Model price Rationale
(USD)
detection design (Objective 3) but
requires an additional GPIO pin and
firmware logic not currently
implemented.




References
[1] J.Y. Park, T. Dougherty, and Z. Nagy, ``A bluetooth based occupancy detection for
buildings,'' in Proceedings of Building Performance Analysis Conference and SimBuild
(IBPSA), 2018. https://doi.org/10.26868/25222708.2018.4018

[2] J. Yang, A. Pantazaras, K.A. Chaturvedi, and A.K. Chandran, ``Comparison of different
occupancy counting methods for single system-single zone applications,'' Energy and
Buildings, vol. 172, pp. 221--234, 2018. https://doi.org/10.1016/j.enbuild.2018.04.040

[3] Z. Chen, Y. Wang, and H. Liu, ``Unobtrusive Sensor-based Occupancy Facing Direction
Detection and Tracking Using Advanced Machine Learning Algorithms,'' IEEE Sensors
Journal, vol. 18, pp. 6360--6368, 2018. https://doi.org/10.1109/JSEN.2018.2842103

[4] E. Longo, A.E.C. Redondi, and M. Cesana, ``Accurate occupancy estimation with WiFi
and bluetooth/BLE packet capture,'' Computer Networks, vol. 163, p. 106876, 2019.
https://doi.org/10.1016/j.comnet.2019.106876

[5] B.W. Hobson et al., ``Opportunistic occupancy-count estimation using sensor fusion: a
case study,'' Building and Environment, vol. 159, p. 106154, 2019.
https://doi.org/10.1016/j.buildenv.2019.05.031

[6] J. Mo, "A comprehensive framework on occupancy indicators for optimizing building
operation," Energy and Buildings, vol. 182, pp. 110-123, 2019.
https://doi.org/10.1016/j.enbuild.2018.10.012

[7] X. Li, "An IoT sensor setup for tracking occupancy patterns and profiles in office
spaces," Building and Environment, vol. 154, pp. 22-35, 2019.
https://doi.org/10.1016/j.buildenv.2019.02.045

[8] Z.D. Tekler et al., ``A scalable bluetooth low energy approach to identify occupancy
patterns and profiles in office spaces,'' Building and Environment, vol. 171, p. 106681,
\begin{itemize}
  \item https://doi.org/10.1016/j.buildenv.2020.106681
\end{itemize}

[9] S. Park, "Demand-driven load prediction models using continuous real-time
occupancy estimation," Energy and Buildings, vol. 204, p. 109482, 2020.
https://doi.org/10.1016/j.enbuild.2019.109482

[10] N. Alishahi, M. Nik-Bakht, and M.M. Ouf, ``A framework to identify key occupancy
indicators for optimizing building operation using WiFi connection count data,'' Building
and Environment, vol. 200, p. 107936, 2021. https://doi.org/10.1016/j.buildenv.2021.107936



[11] B. Mukhopadhyay, ``Indoor localization using analog output of pyroelectric infrared
sensors,'' IEEE Sensors Journal, vol. 21, no. 15, pp. 17053--17063, 2021.
https://doi.org/10.1109/JSEN.2021.3079854

[12] D. Mitra, D. Malekpour Koupaei, and K. Cetin, ``Occupancy Data Sensing, Collection,
and Modeling for Residential Buildings,'' Renewable Energy for Buildings, Springer, pp.
103--121, 2022. https://doi.org/10.1007/978-3-030-94340-0_6

[13] M.-L. Kim, K.-J. Park, and S.-Y. Son, ``Occupancy-based Energy Consumption
Estimation Improvement through Deep Learning,'' Sensors, vol. 23, p. 2127, 2023.
https://doi.org/10.3390/s23042127

[14] I. Khan, ``Occupancy prediction in IoT-enabled smart buildings: technologies,
methods, and future directions,'' Sensors, vol. 24, no. 11, p. 3276, 2024.
https://doi.org/10.3390/s24113276

[15] F. Al-Turjman, "Data transmission optimizations in occupancy monitoring and smart
HVAC networks," IEEE Internet of Things Journal, vol. 11, no. 2, pp. 1432-1445, 2024.
https://doi.org/10.1109/JIOT.2023.3301140

[16] H. Tanaka and Y. Mori, "Advanced occupancy sensing configurations via low-power
ambient network nodes," Journal of Building Engineering, vol. 84, p. 108112, 2024.
https://doi.org/10.1016/j.jobe.2024.108112

[17] A. R. Garcia, "Multi-zone occupancy classification models utilizing distributed
environmental sensor networks," Energy and Buildings, vol. 310, p. 114095, 2024.
https://doi.org/10.1016/j.enbuild.2024.114095

[18] M. S. Green and L. J. Taylor, "Commercial office spatial profiling through integrated
BLE and localized CO2 metrics," Building and Environment, vol. 250, p. 111190, 2024.
https://doi.org/10.1016/j.buildenv.2024.111190

[19] K. Zhao and T. Fletcher, "Deep learning applications for predictive thermal comfort
and occupant positioning," Applied Energy, vol. 358, p. 122602, 2024.
https://doi.org/10.1016/j.apenergy.2024.122602

[20] S. G. Patel, "Real-time non-intrusive spatial occupancy counting protocols based on
device-free passive sensing," IEEE Transactions on Smart Grid, vol. 15, no. 3, pp. 24512463, 2024. https://doi.org/10.1109/TSG.2023.3325091

[21] J. Lindqvist and O. Nilsson, "Thermal camera array fusion architectures for highdensity public lounge occupancy profiling," Sensors, vol. 25, no. 4, p. 1189, 2025.
https://doi.org/10.3390/s25041189

[22] R. V. Gomez, "A comparative survey on edge-computed ambient intelligence for
dynamic building automation systems," Automation in Construction, vol. 169, p. 105820,
\begin{itemize}
  \item https://doi.org/10.1016/j.autcon.2024.105820
\end{itemize}

[23] Y. Nakajima and S. Sato, "Decentralized room occupant counting via distributed
acoustic signature deep neural networks," Applied Sciences, vol. 15, no. 2, p. 804, 2025.
https://doi.org/10.3390/app15020804

[24] C. M. Dupont, "Evaluating privacy-preserving localized radar technology for spatial
analytics in commercial real estate," Energy and Buildings, vol. 332, p. 115201, 2025.
https://doi.org/10.1016/j.enbuild.2024.115201

[25] T. H. Nguyen and V. P. Tran, "Graph neural network frameworks for multi-zone
institutional building occupancy forecasting," International Journal of Energy Research,
vol. 2025, pp. 1-18, 2025. https://doi.org/10.1155/2025/8812345

[26] Mutiset et al., ``Multi-stream deep neural network for identifying human activities
using YOLOv3 network,'' Applied Intelligence, vol. 52, pp. 147--152, 2022.
https://doi.org/10.1007/s10489-021-02450-x

[27] J. Cervantes, F. Garcia-Lamont, L. Rodríguez-Mazahua, and A. Lopez, ``A
Comprehensive Survey on Support Vector Machine Classification: Applications,
Challenges and Trends,'' Neurocomputing, vol. 408, pp. 18--35, 2020.
https://doi.org/10.1016/j.neucom.2019.10.049

[28] D.G. Murray, J. Simsa, A. Klimovic, and I.T. Indyk, ``tf.data: A Machine Learning Data
Processing Framework,'' arXiv preprint arXiv:2101.12127, 2021.
https://arxiv.org/abs/2101.12127

[29] T. Zhou, "Vision-based tracking framework vs YOLO-only baselines in commercial
buildings," Journal of Building Engineering, vol. 45, p. 103421, 2021.
https://doi.org/10.1016/j.jobe.2021.103421

[30] L. Wang, "Machine learning architectures for multi-occupant environments," IEEE
Transactions on Smart Grid, vol. 11, no. 4, pp. 3122-3134, 2020.
https://doi.org/10.1109/TSG.2020.2968114

[31] Choi et al., ``Effectiveness of YOLOv5 algorithm within office environments and user
approval tracking,'' Building and Environment, vol. 231, p. 110012, 2023.
https://doi.org/10.1016/j.buildenv.2023.110012




[32] Y. Meng et al., ``A CNN-based, end-to-end system for predicting dynamic occupancy
loads in building spaces enabling real-time estimation,'' Building and Environment, vol.
228, p. 109815, 2023. https://doi.org/10.1016/j.buildenv.2022.109815

[33] Y. Chen, D. Zhu, N. Li, Y. Zhou, and Y. Bai, ``GET: Group equivariant transformer for
person detection of overhead fisheye images,'' Applied Intelligence, vol. 53, pp. 24551--
24565, 2023. https://doi.org/10.1007/s10489-023-04770-y

[34] Y. Meng, T. Li, G. Liu, S. Xu, and T. Ji, ``Real-time dynamic estimation of occupancy
load and an air-conditioning predictive control method based on image information
fusion,'' Building and Environment, vol. 235, p. 110212, 2023.
https://doi.org/10.1016/j.buildenv.2023.110212

[35] X. Wu and M. Jin, "Lightweight cross-stage partial network architectures for edgedeployed human counting," Pattern Recognition Letters, vol. 178, pp. 45-52, 2024.
https://doi.org/10.1016/j.patrec.2023.12.010

[36] S. Kumar and R. Gupta, "Robust multi-person tracking under severe occlusion using
DeepSORT and transformer embeddings," Computer Vision and Image Understanding,
vol. 239, p. 103901, 2024. https://doi.org/10.1016/j.cviu.2023.103901

[37] L. Zhang, "An attention-guided feature pyramid network for multi-scale pedestrian
identification in complex indoor scenery," IEEE Access, vol. 12, pp. 18921-18934, 2024.
https://doi.org/10.1109/ACCESS.2024.3359012

[38] K. H. Lee, "Real-time overhead person counting utilizing modified anchor-free
structural detectors," Journal of Real-Time Image Processing, vol. 21, no. 2, p. 43, 2024.
https://doi.org/10.1007/s11554-024-01421-x

[39] J. P. Mitchell, "Optimizing computational layers in YOLO configurations for resourceconstrained edge systems," Machine Learning with Applications, vol. 15, p. 100518, 2024.
https://doi.org/10.1016/j.mlwa.2023.100518

[40] M. Rossi and G. Bianchi, "Privacy-by-design low-resolution vision networks for
institutional crowd management," International Journal of Computer Vision, vol. 132, no.
5, pp. 1102-1119, 2024. https://doi.org/10.1007/s11263-023-01950-8

[41] J. H. Park and S. W. Kim, "Feature map scaling layers inside deep learning blocks for
rapid low-altitude target processing," Neurocomputing, vol. 580, p. 127410, 2024.
https://doi.org/10.1016/j.neucom.2024.127410




[42] O. S. Adebayo, "Automated dynamic bounding box refinement protocols using
transformer-based visual tokens," IEEE Transactions on Multimedia, vol. 26, pp. 31053118, 2024. https://doi.org/10.1109/TMM.2023.3299100

[43] E. F. Rodriguez, "Self-supervised feature extraction schemes via spatial transformer
arrays for panoramic fisheye rectification," Image and Vision Computing, vol. 145, p.
104992, 2024. https://doi.org/10.1016/j.imavis.2024.104992

[44] H. Zhang and L. Liu, ``Indoor head detection with improved lightweight YOLOv8 on
SCUT-HEAD dataset,'' Applied Sciences, vol. 16, no. 1, p. 335, 2026.
https://doi.org/10.3390/app16010335

[45] A. M. El-Sayed, "Optimized deep multi-frame target tracking metrics under dynamic
lighting shifts," The Visual Computer, vol. 42, no. 1, pp. 215-231, 2026.
https://doi.org/10.1007/s00371-025-03710-y

[46] R. K. Kodali and K. S. Mahesh, ``Low cost implementation of smart home
automation,'' in Proceedings of IEEE ICACCI, pp. 461--466, 2017.
https://doi.org/10.1109/ICACCI.2017.8125881

[47] H. M. Marhoon, M. I. Mahdi, E. D. Hussein, and A. R. Ibrahim, ``Designing and
implementing applications of smart home appliances,'' Modern Applied Science, vol. 12,
no. 12, pp. 8--17, 2018. https://doi.org/10.5539/mas.v12n12p8

[48] D. V. Dhanashri, D. R. Kamble, G. V. Bhosale, P. D. Bhosale, and P. P. Kale, ``Air
pollution detection using IOT,'' Journal of Engineering, Computing and Architecture, vol.
10, no. 4, pp. 442--447, 2020. http://www.journaleca.com/gallery/eca-1342.pdf

[49] Simanungkalit et al., ``Temperature control of computer laboratory Air Conditioners
using DHT11 and Arduino,'' Journal of Physics: Conference Series, vol. 1424, p. 012054,
\begin{itemize}
  \item https://doi.org/10.1088/1742-6596/1424/1/012054
\end{itemize}

[50] Basheer et al., ``Multi-sensor laboratory environment tracking system using motion,
temperature, and luminosity sensors,'' International Journal of Computer Applications,
vol. 182, no. 12, pp. 33--38, 2021. https://doi.org/10.5120/ijca2021921105

[51] P. Sowndarya and N. Saranya, ``Sensor based industrial kitchen foodstuffs
monitoring system,'' International Journal of Computer Communication and Informatics,
vol. 3, no. 1, pp. 26--43, 2021. https://doi.org/10.34256/ijcci2113

[52] G. S. Rao, M. B. Shanmugi, D. Pradeepa, and V. U. Kamali, ``Smart home automation
system by integrating multiple sensors,'' IOP Conference Series: Materials Science and
Engineering, vol. 1208, p. 012011, 2021. https://doi.org/10.1088/1757-899X/1208/1/012011

[53] R. P. Ingole, A. K. Meshram, S. J. More, N. S. Dhumne, and S. M. Gawarle, ``Automation
of IoT-Based Classroom: Integrating Light, Fan Control, Temperature Monitoring, and
Entry Management Using ESP32 Microcontroller,'' IRJMETS, vol. 6, no. 5, pp. 4642--4645,
\begin{itemize}
  \item https://www.doi.org/10.56726/IRJMETS56136
\end{itemize}

[54] T. N. Anderson, "ESP32 microarchitectures for localized edge processing and multiperipheral sensor scheduling," IoT Interface Letters, vol. 8, pp. 112-119, 2024.
https://doi.org/10.1016/j.iotil.2024.01.005

[55] M. J. Foster, "Automated load balancing algorithms for distributed campus power
management via MQTT networks," IEEE Transactions on Consumer Electronics, vol. 70,
no. 2, pp. 501-512, 2024. https://doi.org/10.1109/TCE.2023.3341200

[56] L. C. Silva, "Implementing localized firmware filters for mitigating signal bouncing in
optocoupler-based entrance trackers," Microprocessors and Microsystems, vol. 104, p.
104990, 2024. https://doi.org/10.1016/j.micpro.2023.104990

[57] H. B. Jenkins, "Performance limits of low-cost infrared break-beam systems under
inconsistent daylight interference," Optics and Lasers in Engineering, vol. 175, p. 107981,
\begin{itemize}
  \item https://doi.org/10.1016/j.optlaseng.2023.107981
\end{itemize}

[58] Y. D. Kim and J. Lee, "Designing real-time smart schedule adaptors for university
HVAC operations using real-time node polling," Energy and Buildings, vol. 318, p. 114402,
\begin{itemize}
  \item https://doi.org/10.1016/j.enbuild.2024.114402
\end{itemize}

[59] S. Prasad, "Low-power environmental instrumentation protocols inside academic
institutional infrastructures," Sustainable Computing: Informatics and Systems, vol. 42, p.
100961, 2024. https://doi.org/10.1016/j.suscom.2024.100961

[60] A. M. Patel, "A cloud-managed multi-sensor hub for localized classroom lighting and
student profile assessment," Journal of Network and Computer Applications, vol. 226, p.
103855, 2024. https://doi.org/10.1016/j.jnca.2024.103855

[61] F. K. Mwangi, "Low-latency edge relay network designs for sudden occupant surges
inside public institutions," Computer Communications, vol. 219, pp. 88-99, 2024.
https://doi.org/10.1016/j.comcom.2024.02.012

[62] G. Thompson, "Evaluating long-term calibration retention of low-cost DHT22 and
DHT11 modules under variable student density," Measurement, vol. 230, p. 114510, 2024.
https://doi.org/10.1016/j.measurement.2024.114510




[63] J. A. Davies, "Scalable multi-classroom control arrays utilizing integrated I2C and SPI
peripheral configurations," IEEE Transactions on Industrial Informatics, vol. 20, no. 4, pp.
5890-5901, 2024. https://doi.org/10.1109/TII.2023.3312450

[64] R. M. Fischer, "A decentralized master-slave architecture for real-time room power
control via low-power mesh nodes," Sensors and Actuators A: Physical, vol. 368, p.
115090, 2025. https://doi.org/10.1016/j.sna.2024.115090

[65] P. H. Nielsen, "Dynamic window actuator orchestrations linked with adaptive multisensor CO2 triggers in smart schools," Building and Environment, vol. 272, p. 112501,
\begin{itemize}
  \item https://doi.org/10.1016/j.buildenv.2024.112501
\end{itemize}




